% ==========================================================
% 不在首位且两个对象不相邻
% 难度：★★ 中档
% ==========================================================

\newcommand{\qProbPermNotFirstNotAdjStem}{%
“礼、乐、射、御、书、数”六次讲座各安排一次，要求：“射”不在第一次；“数”和“乐”不相邻。求不同讲座次序数。
}

\newcommand{\qProbPermNotFirstNotAdjAnswer}{%
$408$
}

\newcommand{\qProbPermNotFirstNotAdjSolution}{%
先只满足“射不在第一次”：\\
“射”可以安排在第 $2$ 至第 $6$ 个位置（共 $5$ 个）；\\
其余五门讲座任意排列。\\
方法数为：$5\cdot5!=600$ 种。

下面从中减去“数”和“乐”相邻的情况：\\
把“数”和“乐”捆绑成一个整体，内部有 $2!$ 种。\\
捆绑后共有五个单位（包含“射”）。\\
要求“射”不在第一个单位位置：
\begin{itemize}
    \item “射”可选择后四个位置，有 $4$ 种；
    \item 其余四个单位排列，有 $4!$ 种；
    \item “数乐”整体内部有 $2$ 种顺序。
\end{itemize}
不合要求的方法数为：$4\cdot4!\cdot2=192$ 种。

因此合法的不同讲座次序数为：
\[ N = 600 - 192 = 408 \]
}

\newcommand{\qProbPermNotFirstNotAdjTeachingNotes}{%
\textbf{【避坑与边界说明】}
\begin{itemize}
    \item “数”和“乐”捆绑后仍有两种内部次序。
    \item 捆绑后共有五个单位，而不是四个。
    \item 减去坏情况时，不能只减去“数乐相邻”。\\
    必须同步保留“射不在第一位”的条件。
    \item “先满足一个，再排除另一个不满足”，比逐个讲座插空更加明了。
\end{itemize}
}
